\documentclass[11pt]{article}
\usepackage{amsmath,amsfonts}
\usepackage{graphicx}
%\usepackage{xcolor}
%\definecolor{gray}{rgb}{.75,.75,.75}
%\usepackage[landscape]{geometry}
%\usepackage{hyperref}
%\usepackage[bookmarks=true, pdfborder={0 0 0}, urlbordercolor=gray]{hyperref}

\renewcommand{\thefootnote}{\fnsymbol{footnote}}

\addtolength{\topmargin}{-.4in}
\addtolength{\textheight}{.7in}
\addtolength{\oddsidemargin}{-.2in}
\addtolength{\textwidth}{.4in}
\pagestyle{empty}

\newcommand{\ds}{\displaystyle}
\newcommand{\ts}{\textstyle}

\newcommand{\R}{\mathbb{R}}
\newcommand{\tZ}{\tilde{\Z}}
\newcommand{\sgn}{\mathrm{sgn}}

\newcommand{\Sym}{\mathrm{Sym}}
\newcommand{\tu}{\textunderscore}

\renewcommand{\circle}[1]{{\bigcirc\hspace{-.75em}#1}}
\renewcommand{\epsilon}{\varepsilon}
\newcommand{\st}{\mbox{\footnotesize$\sqrt3$}}

\newcommand{\ul}{\underline}

\begin{document}

\footnote[0]{Notes by Peter Magyar \ {\tt magyar@math.msu.edu}}
\vspace{-1em}

\centerline{\bf Math 132\hspace{1.5em} \hfill{\Large\bf Area Between Curves}\hfill Stewart \S5.1}
\vspace{1em}

\noindent
{\bf Region between two parabolas.} We have seen that geometrically, the integral $\int_a^b f(x)\,dx$ computes the area between a curve $y=f(x)$ and an interval $x\in[a,b]$ on the $x$-axis (with area below the axis counted negatively). In Calculus II, we will show the versatility of the integral to compute all kinds of areas, lengths, volumes: almost any measure of size for a geometric object.

In this section, we compute more general areas:
those between two given curves $y=f(x)$ and $y=g(x)$,
usually with no boundary on the $x$-axis.
\\[.5em]
{\sc example:} Consider the region 
with top boundary $y=f(x) =x^2{+}x{+}1$, 
bottom boundary $y=g(x)=2x^2{-}1$, 
left boundary the $y$-axis $x\,{=}\,0$,
right boundary~$x\,{=}\,1$:\footnote{
We specify the region as the set  of 
all points $(x,y)$ which satisfy the given
conditions.
}
$$\begin{array}{rcl}
R &=& 
\{\,(x,y) \ \text{with}\ g(x)\leq y\leq f(x)
\text{ and } x\in[0,1]\,\}.
\\[.5em]
&=& 
\{\,(x,y) \ \text{with}\ 2x^2{-}1\leq y\leq x^2{+}x{+}1
\text{ and } 0\leq x\leq 1\,\}.
\end{array}$$
\\[0em]
\centerline{
\includegraphics[width=3.8in]
{5-1-Parabolas-1.png}
%\hspace{.5in}
%\includegraphics[width=2in]
%{5-1-Vert-Slices.png}
}
\vspace{-.5em}

\noindent
Here $y=g(x)=2x^2{-}1$ is a standard parabola shifted downward, with minimum point $x=0$.
The curve $y=f(x)=x^2{+}x{+}1$ is roughly like its leading term $y=x^2$, 
a parabola opening upward; 
its minimum point satisfies $(x^2{+}x{+}1)' = 2x{+}1 = 0$, i.e.~$x=-\frac12$.

To compute the area of $R$, we use the same geometric-numerical strategy as for the region under a single curve: split $R$ into $n$ thin vertical slices of width $\Delta x =\frac1n$, each approximately a rectangle;
then add up the rectangle areas and take the limit
as $n$ becomes larger and larger.
In the interval $x\in[0,1]$, we take sample points $x_1,\ldots, x_n$, one in each $\Delta x$ increment.
The slice at position $x_i$ has height equal to the ceiling minus the floor,
$f(x_i)-g(x_i)$, so:
$$
\text{area of slice} \ \approx\ \text{(height)$\times$(width)} \ =\  (f(x_i)-g(x_i))\,\Delta x,
$$
and the total area is:
$$
A_{[0,1]} \ = \ \lim_{n\to\infty} \sum_{i=1}^n (f(x_i)-g(x_i))\,\Delta x
\ =\ \int_0^1\! (f(x)-g(x))\,dx
$$
$$ \ =\
\int_0^1 (x^2{+}x{+}1)-(2x^2{-}1)\,dx \ =\ \left[\vphantom{\sum}2x{+}\tfrac12x^2{-}\tfrac13x^3\right]_{x=0}^{x=1} \ =\ \tfrac{13}6\,.
$$
\\[.5em]
{\sc example:}
Next, consider the region between the same curves
$y=f(x) =x^2{+}x{+}1$ and $y=g(x)=2x^2{-}1$,
but above the interval $x\in[1,3]$. 
To picture the region without a calculator,
we determine the intersection points where the  curves cross:
$$
f(x) = g(x)
\quad\Longleftrightarrow\quad
x^2{+}x{+}1 = 2x^2{-}1
\quad\Longleftrightarrow\quad
$$
$$
x^2{-}x{-}2 = 0
\quad\Longleftrightarrow\quad
x = \tfrac{1\pm\sqrt{1^2{-}4(1)(-2)}}{2(1)}= \tfrac{1\pm3}2
= -1\text{ or }2
$$
by the Quadratic Formula. Only $x=2$ is relevant for our region above $x\in[1,3]$.
At $x=1$ we have $g(1)<f(1)$,
so to the left of $x=2$, our region is defined by $g(x)\leq y\leq f(x)$.
At $x=3$, we have $f(3)<g(3)$, 
so to the right of $x=2$, it is $f(x)\leq y\leq g(x)$: 
\\[.5em]
\centerline{
\includegraphics[width=3.8in]
{5-1-Parabolas-2.png}
}
\vspace{-.5em}

\noindent
Repeating our previous area formula for the two parts of our region gives:
$$
A_{[1,3]} \ = \  A_{[1,2]} + A_{[2,3]}
\ = \ 
\int_1^2 (f(x)-g(x))\,dx + \int_2^3 (g(x)-f(x))\,dx
$$
$$
\ =\ \int_1^2 (2{+}x{-}x^2) \,dx \,+ \int_2^3 (-2{-}x{+}x^2)\,dx
\ =\
\tfrac76 + \tfrac{11}6 \ =\ 3.
$$
\newpage

\noindent
{\sc example:} Finally, we consider the same curves
$y=f(x) =x^2{+}x{+}1$ and $y=g(x)=2x^2{-}1$, 
but we take the entire finite region between
them: $$R'=\{(x,y)\text{ with } g(x)\leq y\leq f(x)\}.$$
\\[-1em]
\centerline{
\includegraphics[width=3.7in]
{5-1-Parabolas-3.png}
}
\vspace{-.5em}

\noindent
Here the top boundary is $y=x^2+x+1$ and the bottom boundary is $y=2x^2-1$,
but we have not specified an $x$ interval. However, we have already computed
the intersection points $x=-1$ and $x=2$, and the curves
do not enclose any finite regions beyond these points. Thus:
$$
A_{[-1,2]} \ =\ \int_{-1}^2 (f(x)-g(x))\,dx = \tfrac92\,.
$$ 

\noindent
We can generalize the above examples in:
\begin{quote}
{\it Theorem:} The area of the region enclosed between $f(x)$ and $g(x)$ for $x\in[a,b]$
is:
$\ds
A\ =\ 
\int_a^b \left|\vphantom{\tfrac AA}f(x)-g(x)\right|\,dx\,.
$
\end{quote}

\noindent
The absolute value signs ensure we take the integral of top minus bottom, regardless of 
which is which.
In practice, we must find the intersection points where $f(x)=g(x)$,
which split the integral into intervals where $g(x)\leq f(x)$ versus
$f(x)\leq g(x)$.
\\[1.5em]
{\bf Integrating with respect to $\mathbf{y}$.}
Consider the region:
$$
R \ =\ \{(x,y) \text{ with }  y^2 \leq x  \leq y{+}1\}\,.
$$
Here the boudary curves are naturally graphs in which $y$ is the independent variable:
the right boundary is the line $x=f(y)=y{+}1$;
and the left boundary is $x = g(y) = y^2$, a parabola opening to the right.

\addtolength{\textheight}{8em}

\newpage

\addtolength{\topmargin}{-4em}


\centerline{
\includegraphics[width=2in]
{5-1-Y-Function.png} 
}
\vspace{-0em}

\noindent
Understand: it is merely by habit that we consider $y$ as a function of $x$.
We can make $x$ a function of $y$ instead if it is more convenient,
and the same formulas will work if we switch the roles of $x$ and $y$. 
Thus, we find the intersection points:
$y{+}1 = y^2$ when $y=\frac{1{\pm}\sqrt5}2$ by the Quadratic Formula.
The area as:
$$
A  \ =\ \int_{\frac{1{-}\sqrt5}2}^{\frac{1{+}\sqrt5}2} (y{+}1) - (y^2)\,dy
\ = \ \left[\vphantom{\int}\tfrac12 y^2 + y - \tfrac13y^3 \right]_{y=\frac{1{-}\sqrt5}2}^{y=\frac{1{+}\sqrt5}2}
\ =\  \tfrac56\sqrt5.
$$
Here $((y{+}1)- (y^2))\,dy$ represents the area of the {\it horizontal} slice of the region at height $y$,
with thickness $dy$.

To check this, we re-do it from our usual perspective, using $x$ as the independent variable. This makes it more complicated, since we must consider the 
region as having three boundary graphs: 
upper boundary $y=\sqrt x$,  lower right boundary $y=x{-}1$,
and lower left boundary $y=-\sqrt x$. The intersection points are:
\\[-1.5em]
\begin{itemize}
\item Between $y=\sqrt x$ and $y=x{-}1$: \   $x = \frac{3{+}\sqrt5}{2}$ \ (upper right corner)
\\[-1.5em]
\item Between $y=-\sqrt x$ and $y=x{-}1$: \   $x = \frac{3{-}\sqrt5}{2}$ \ (lower middle corner)
\\[-1.5em]
\item Between $y=\sqrt x$ and $y=-\sqrt x$: \  $x=0$ \ (left end)
\end{itemize}
\vspace{-0.3em}

\noindent
These split the region into left and right parts:
\\[1em]
\centerline{
\includegraphics[width=2in]
{5-1-X-Function.png} 
}
\vspace{-1em}

\noindent
The area is:
\\[-1.0em]
$$
A \ =\ \int_{0}^{\frac{3{-}\sqrt5}{2}} (\sqrt x) - (-\sqrt x)\,dx
+ \int_{\frac{3{-}\sqrt5}{2}}^{\frac{3{+}\sqrt5}{2}} (\sqrt x) - (x{-}1)\,dx\,,
$$
which after much algebra gives the same answer as before: $\tfrac56\sqrt5$.



\end{document}

%%%%%%%%%%%  Picture  %%%%%%%%%%
\\[1em]
\centerline{
%\includegraphics[width=2in]
{1-8-Discont-iv.png}
\qquad \qquad 
%\includegraphics[width=2in]
{1-8-Discont-v.png} 
}
\vspace{1em}

\noindent

%%%%%%%%%%%  Table  %%%%%%%%%%
$$\begin{array}{|c|c|}
A & B
\\ \hline && \\[-1.1em]
C & D
\end{array}$$

%%%%%%%%%%% Outline %%%%%%%%%%%%
\begin{enumerate}

\item  \begin{itemize}
\item
\item 
\item 
\end{itemize}

\end{enumerate}

\\[1em]
\end{itemize}

